\documentclass[12pt]{article}
\usepackage[T1]{fontenc}
\usepackage[utf8]{inputenc}
\usepackage{lmodern}
\usepackage{textcomp}
\usepackage{enumitem}
\usepackage{geometry}
\usepackage{graphicx}
\usepackage{amsmath, amssymb}
\geometry{margin=1in}
\begin{document}
\begin{center}
Biology - Undergraduate Level Assessment 20 \
Total Marks: 40 \
Answer all questions.
\end{center}

\section*{Multiple Choice (10 marks)}
\begin{enumerate}[label=\arabic*.]
\item If a random variable X is normally distributed with a mean of 118 and a standard deviation of 11, what Z-scores correspond to raw scores of 115, 134, and 99?
\begin{enumerate}[label=(\alph*)]
    \item -.27, 1.45, 1.73
    \item -.27, -1.45, -1.73
    \item -2.18, 0.15, -1.82
    \item 0, 1.5, -2
    \item .27, 1.45, 1.73
\end{enumerate}
\item This biome contains plants whose roots cannot go deep due to the presence of a permafrost.
\begin{enumerate}[label=(\alph*)]
    \item Savanna
    \item Taiga
    \item Tundra
    \item Chaparral
    \item Alpine
\end{enumerate}
\item What are the major forces of evolution?
\begin{enumerate}[label=(\alph*)]
    \item reproduction, migration, mutation, natural selection
    \item mutation, genetic drift, migration, genetic hitchhiking
    \item mutation, genetic drift, migration, random mating
    \item mutation, gene flow, migration, natural selection
    \item adaptation, genetic drift, migration, natural selection
\end{enumerate}
\item Which group is composed entirely of individuals who maintained that species are fixed (i.e., unchanging)?
\begin{enumerate}[label=(\alph*)]
    \item Aristotle, Darwin, and Lamarck
    \item Aristotle, Lyell, and Darwin
    \item Aristotle, Linnaeus, and Cuvier
    \item Aristotle, Cuvier, and Lamarck
    \item Linnaeus, Darwin, and Lamarck
\end{enumerate}
\item Fungi participate in each of the following EXCEPT
\begin{enumerate}[label=(\alph*)]
    \item photosynthesis to produce glucose
    \item fermentation to produce alcohol
    \item association with humans to produce ringworm
    \item association with the roots of plants to form mycorrhizae
\end{enumerate}
\end{enumerate}

\section*{True / False (10 marks)}
\textit{Answer True or False}
\begin{enumerate}[label=\arabic*.]
\setcounter{enumi}{5}
\item True or False: Gene flow between populations increases the likelihood of speciation by enhancing genetic divergence.
\item True or False: Reducing the carrying capacity of the environment can effectively lower the K value and help control mosquito populations.
\item True or False: The extra Y chromosome is linked to an increase in artistic skills.
\item True or False: Commensalism is a form of interspecies interaction where one species benefits while the other is neither helped nor harmed.
\item True or False: Hydrolysis is involved in the conversion of starch to simple sugars.
\end{enumerate}

\section*{Long Form Response (20 marks)}
\textit{Answer each question with a clear and complete explanation.}
\begin{enumerate}[label=\arabic*.]
\setcounter{enumi}{10}
\item Discuss the methods used to estimate the age of a rock, focusing on the role of radioactive elements and fossil content in determining geological timelines.
\item Discuss the concepts of incomplete dominance and epistasis in genetics, analyzing how they differ in terms of the interactions between alleles and genes.
\end{enumerate}

\vfill
\noindent\textit{End of Paper}
\end{document}